Данная работа рассматривала генерацию юмора как задачу рассуждения в неверифицируемом домене. Центральное утверждение состояло в том, что обучение юмору должно отходить от имитации одного референса и переходить к распределению вероятностной массы по множеству референсов при слабых запросах. Обзор литературы показал, почему такой сдвиг необходим: SFT склонен к имитации и запоминанию, preference-based RL нестабилен в юморе из-за шумных оценок, а современные многошаговые системы генерации юмора повышают качество ценой жестко заданных структур рассуждения.

Работа внесла два связанных вклада. На методологическом уровне она сформулировала multi-reference расширение verifier-free обучения с подкреплением, в котором reward задается полной вероятностной массой, назначенной множеству приемлемых референсов, а не одному золотому ответу. На уровне реализации она построила и проанализировала пилотный пайплайн, который преобразует большой корпус шуток в prompt-reference наборы через embeddings, извлечение ключевых слов, vector retrieval, deduplication и разбиение данных.

Эмпирический анализ пилота показал, что идея является практически осуществимой. В текущем локальном состоянии репозитория уже собран корпус из 664{,}048 шуток, пилотное подмножество было embedded и обработано, а итоговый workflow производит 5{,}058 дедуплицированных prompt-reference строк со средним числом референсов 4.07 на prompt и медианой пять. Эти статистики показывают, что ключевое структурное требование MRVF learning, а именно наличие prompts с несколькими приемлемыми референсами, может быть удовлетворено на практике.

В то же время работа остается предварительной в корректном смысле. Процедура построения prompts все еще приближенная, downstream benchmark artifacts еще не закоммичены, а сам MRVF training loop пока не реализован. Поэтому данную работу следует читать как основу дипломного исследования, а не как окончательный эмпирический вердикт об обучении с подкреплением для юмора.

Главный вывод состоит в том, что проект уже перешел границу от спекулятивного отчета к защищаемой исследовательской программе. Представление данных, программный пайплайн и формальная objective согласованы друг с другом. Следующий шаг состоит в реализации этапа обучения и проверке того, сможет ли построенный multi-reference датасет действительно вызывать более качественное рассуждение и более сильный юмор, чем стандартные post-training baselines.
